\documentclass[12pt]{article}
\usepackage[utf8]{inputenc}
\usepackage[T1]{fontenc}
\usepackage{hyperref}
\title{Pesse Language Resources for K-12 Curriculum Development}
\author{Compiled by Claude Sonnet 5.5 (claude-sonnet-5-5) for Andr\'as Kornai}
\date{}
\begin{document}
\maketitle

\section{Introduction}
Pesse (ISO 639-3 pze) is an Afro-Asiatic (Chadic) language of Bauchi State, Nigeria. Search results indicate that speakers call themselves Pesse, while the government and neighbouring groups call them Polchi/Polci and the Hausa call them Polsawa. Settlements are reported in the Polchi District (Dass LGA) and Toro District (Toro LGA). Speaker estimates vary (about 37,000 to 41,000 on Joshua Project pages; about 50,000 for Polci on Verbix), and the relation between Pesse and the Polci cluster is not fully clear from the sources found. Very little educational material was found.

\section{Primary Educational Resources}
No primers, textbooks or school programmes for Pesse were found in the searches. Polci is reported to be giving way to Hausa, and speakers are largely bilingual in Hausa.

\section{Linguistic and Reference Materials}
\begin{itemize}
\item Wikipedia, Polci language (general overview, classification in the Barawa cluster of West Chadic): \url{https://en.wikipedia.org/wiki/Polci_language}
\item Endangered Languages Project entry for Polci (alternative names, classification): \url{https://endangeredlanguages.com/lang/4975}
\item Verbix, Endangered Languages: Polci: \url{https://docs.verbix.com/EndangeredLanguages/Polci}
\end{itemize}

\section{Religious and Community Resources}
\begin{itemize}
\item Joshua Project, Pesse in Nigeria (people group profile): \url{https://joshuaproject.net/people_groups/14430}
\end{itemize}
No scripture translations or community organisations were identified in the search results.

\section{Digital and Audio Resources}
\begin{itemize}
\item Joshua Project language page for pze: \url{https://joshuaproject.net/languages/pze}
\end{itemize}
No audio, video or corpus resources were found.

\section{Curriculum Development Context}
There is no known Pesse-medium schooling. A curriculum would need to start from basic documentation (orthography, wordlists, recordings) in cooperation with the community, taking account of the dominance of Hausa and the unclear relationship with Polci.

\end{document}